\section{Risk Register}
\label{app:risk}

The following table identifies the principal project risks, scored using the standard
likelihood-severity matrix (L\,=\,1, M\,=\,2, H\,=\,3; Score\,=\,Likelihood\,$\times$\,Severity).

\begin{table}[H]
\centering
\caption{Project Risk Register}
\label{tab:risk}
\small
\begin{tabularx}{\textwidth}{p{3.6cm} c c c X}
\toprule
\textbf{Project Risk} &
\multicolumn{2}{c}{\textbf{Likelihood}} &
\multicolumn{2}{c}{\hspace{-1cm}} \\
\cmidrule(lr){2-3}
 & \textbf{L\;M\;H} & \textbf{L\;M\;H} & \textbf{Score} & \textbf{Mitigation Measures} \\
 & \small(Likelihood) & \small(Severity) & \small(L$\times$S) & \\
\midrule
\textbf{Simulation instability:} environment bugs cause non-physical state transitions or incorrect reward computation.
& M & H & 6 &
Unit-test each component (propagation, ADR, Capture Effect) in isolation. Episode step counters and boundary clipping enforce plausible states. \\
\midrule
\textbf{Training divergence:} DQN fails to converge due to reward scale mismatch or exploding gradients.
& M & H & 6 &
Reward normalisation; gradient clipping; target network stabilisation; hyperparameter sweep on learning rate and batch size. \\
\midrule
\textbf{Overfitting to training scenarios:} policy memorises fixed layouts, fails to generalise.
& H & H & 9 &
Domain Randomisation: sensor positions, grid sizes, and counts randomised every episode. Curriculum Learning ensures gradual exposure to difficulty. \\
\midrule
\textbf{Computational resource constraints:} long training runs exceed available GPU time.
& L & M & 2 &
Primary training on personal RTX 3050 Ti; long runs scheduled overnight. Stable-Baselines3 provides efficient GPU-accelerated training. \\
\midrule
\textbf{Sim-to-Real gap:} policy performs well in simulation but fails on a physical UAV.
& M & M & 4 &
High-fidelity models (Two-Ray, log-normal shadowing, EMA-ADR) reduce the gap. Domain Randomisation adds robustness. Acknowledged as limitation; HIL validation proposed as future work. \\
\midrule
\textbf{Time overrun:} evaluation and write-up take longer than planned.
& M & M & 4 &
Modular dissertation structure allows parallel writing. Regular supervisor meetings track progress. Scope adjusted if needed. \\
\bottomrule
\end{tabularx}
\end{table}

\section{Health \& Safety Risk Assessment}
\label{app:hs}

This project is entirely computational and involves no physical experiments, hazardous materials, human participants, or off-campus fieldwork. The assessment below was completed in accordance with the University's online health and safety induction.

\begin{table}[H]
\centering
\caption{Health \& Safety Risk Assessment}
\label{tab:hs_risk}
\small
\begin{tabularx}{\textwidth}{p{2.4cm} p{2.2cm} p{2.0cm} X p{1.1cm} p{0.8cm}}
\toprule
\textbf{Activity} & \textbf{Hazard} & \textbf{Who might be harmed and how} & \textbf{Existing measures to control risk} & \textbf{Risk rating} & \textbf{Result} \\
\midrule
Regular computer use (coding, training, writing)
& Display screen equipment (DSE): eye strain, musculoskeletal discomfort, fatigue
& Student: prolonged screen exposure and poor posture
& Follow DSE guidance: screen at arm's length, top of screen at eye height; take 5--10 min break every hour; use ergonomic seating; adjust screen brightness to ambient lighting.
& Low & A \\
\midrule
Use of electrical equipment (laptop, workstation, monitors)
& Electrical faults: electric shock, fire risk from damaged cables
& Student and others nearby
& Use only PAT-tested equipment; inspect cables before use; do not use damaged cables or adaptors; do not overload sockets; switch off equipment when not in use.
& Low & A \\
\midrule
Lone and off-campus working (home and university library)
& Difficulty obtaining help in an emergency; isolation
& Student
& Carry a charged mobile phone at all times; work in staffed university buildings during core hours where possible; inform a contact of working location; follow University lone-worker policy.
& Low & A \\
\bottomrule
\end{tabularx}
\end{table}

\noindent\textbf{Overall risk level:} Low. No hazardous materials, fieldwork, human participants, or physical experiments are involved. No special permits or ethical approval are required.
